\section*{الفصل \ref{ch:b3:colloids} --- السطوح الفاصلة والمواد الفعالة سطحيًا والغروانيات}

\begin{solution}{exo:b3:colloids:1}
$r = \qty{0.50}{\micro m}$: $\Delta p = 2 \times 0.0720/\num{5.0e-7} = \qty{2.9e5}{Pa}$؛ وفي الداخل، $1.0 + 2.9 = \qty{3.9}{bar}$.
\end{solution}

\begin{solution}{exo:b3:colloids:2}
$\ln(p/p^*) = 2 \times 0.0720 \times \num{1.807e-5}/(\num{1.0e-8} \times 8.314 \times 298.15) = 0.105$؛ $p/p^* = 1.11$.
\end{solution}

\begin{solution}{exo:b3:colloids:3}
$\cos\theta = (40 - 20)/72 = 0.278$، $\theta = \ang{74}$: أقل من $90^\circ$، فالسائل \omterm{def:b3:colloids:wetting}{يبلّل} الصلب
جزئيًا.
\end{solution}

\begin{solution}{exo:b3:colloids:4}
NaCl: $I = \qty{10}{mol.m^{-3}}$، $\kappa^{-1} = 0.96\sqrt{100/10} = \qty{3.0}{nm}$. \ce{MgSO4}: $I = \qty{40}{mol.m^{-3}}$، \qty{1.5}{nm}.
\end{solution}

\begin{solution}{exo:b3:colloids:5}
$\Gamma = \num{12.6e-3}/(2 \times 8.314 \times 298.15) = \qty{2.54e-6}{mol.m^{-2}}$؛ والمساحة $1/(\Gamma N_A) =
\qty{0.65}{nm^2}$ لكل جزيء.
\end{solution}

\begin{solution}{exo:b3:colloids:6}
دون \omterm{def:b3:colloids:micelle}{التركيز المذيلي الحرج} يكون الميل $-12.9$ \unit{mN/m} لكل وحدة من $\log c$ (من $-5.0$ إلى $-4.0$)؛ 
والهضبة عند 33.4. ويبلغها المستقيم عند $\log c = -4.0 + (39.1 - 33.4)/12.9 = -3.56$: \omterm{def:b3:colloids:micelle}{فالتركيز المذيلي الحرج}
$\approx \qty{2.8e-4}{mol/L}$.
\end{solution}

\begin{solution}{exo:b3:colloids:7}
دوديسيل كبريتات الصوديوم: $P = 0.350/(0.60 \times 1.67) = 0.35$، على الحد بين الكرات والأسطوانات
القصيرة: مذيلات كروية قرب \omterm{def:b3:colloids:micelle}{التركيز المذيلي الحرج}. الليبيد: $P = 0.700/(0.65 \times 1.67) = 0.64$:
طبقات ثنائية، ومنه حويصلات.
\end{solution}

\begin{solution}{exo:b3:colloids:8}
\ce{CaCl2}: $60/64 = \qty{0.94}{mM}$ من \ce{Ca^{2+}}؛ \ce{AlCl3}: $60/729 = \qty{0.082}{mM}$ من \ce{Al^{3+}}. أما \ce{Na3PO4}
فيخثّر بأيوناته \ce{Na+} (عند نحو \qty{20}{mM} من الملح)؛ والفوسفات، وهو أيون
مرافق، قد يُمتز بل ويزيد الشحنة السالبة.
\end{solution}

\begin{solution}{exo:b3:colloids:9}
الجزء المحب للماء $6 \times 44.0 + 17.0 = \qty{281.0}{g/mol}$ من \qty{450.0}{g/mol} (الكتل الذرية المعتمدة في الكتاب):
HLB $= 20 \times 281.0/450.0 = 12.5$، أي مستحلِب زيت في ماء.
\end{solution}

\begin{solution}{exo:b3:colloids:10}
زيت القطيرات الصغيرة أكثر ذوبانًا في الطور المتصل (بالعامل
$\exp(2\gamma V_{\mathrm m}/rRT)$، الذي أُسّه أكبر بعشر مرات من أجل \qty{50}{nm} منه من أجل \qty{500}{nm}): فهو
ينتشر نحو القطيرات الكبيرة، التي تنمو بينما تختفي الصغيرة. وهذا
\omterm{def:b3:colloids:ripening}{الإنضاج الأوستفالدي} يُخشّن \omterm{def:b3:colloids:colloid}{المستحلب} حتى لو لم تندمج قطيرتان قط.
\end{solution}

\begin{solution}{exo:b3:colloids:11}
الهلالة غطاء كروي نصف قطره $r/\cos\theta$: فالسائل تحتها مباشرة عند
ضغط أقل بمقدار $2\gamma\cos\theta/r$، ويرتفع العمود حتى يعوّض ذلك $\rho gh$. ومن أجل الماء،
$h = 2 \times 0.072/(997 \times 9.81 \times \num{1.0e-4}) = \qty{0.15}{m}$.
\end{solution}

\begin{solution}{exo:b3:colloids:12}
عند \qty{100}{mM} يكون \omterm{def:b3:electrode-kinetics:double-layer}{طول ديباي} دون \qty{1}{nm} ويتلاشى التنافر قبل
التجاذب: فلا يبقى أي حاجز (أبعد من منحنى \qty{50}{mM}). والبوليمر الحر لا
يعمل عند السطح؛ أما طبقة البوليمر المطعَّمة فتتنافر عند التماس مهما كان
الملح، بسبب الإنتروبيا المفقودة حين تتداخل السلاسل.
\end{solution}

\begin{solution}{pb:b3:colloids:1}
\textbf{1.} الاستبدالات في الشبكة البلورية (\ce{Al^{3+}} مكان \ce{Si^{4+}}، و\ce{Mg^{2+}} مكان \ce{Al^{3+}}) تترك
شحنة سالبة دائمة، تعوّضها كاتيونات قابلة للتبادل.
\textbf{2.} $I = \frac12(1.0 + 1.0) = \qty{1.0}{mM}$.
\textbf{3.} $\kappa^{-1} = \qty{9.6}{nm}$.
\textbf{4.} $a\kappa = 10$: فالطبقة المزدوجة رقيقة مقارنةً بالجسيم، وهذا
شرط تقريب ديرياغين.
\textbf{5.} $v = 2 \times 1650 \times 9.81 \times (\num{1.0e-7})^2/(9 \times \num{0.89e-3}) = \qty{4.0e-8}{m/s}$، أي نحو \qty{3.5}{mm} في اليوم.
\textbf{6.} تتنافر طبقاتها المنتشرة قبل أن يغلب تجاذب فان دير فالس.
\textbf{7.} $V = 2\pi\varepsilon_r\varepsilon_0a\psi^2\ln(1 + \eu^{-\kappa h}) - Aa/12h$.
\textbf{8.} نحو $39\,kT$.
\textbf{9.} جزء من رتبة $\eu^{-39} \approx 10^{-17}$ من التصادمات: أي عمليًا أبدًا.
\textbf{10.} يتلاشى: فكل تصادم يؤدي إلى التماس.
\textbf{11.} \omterm{def:b3:electrode-kinetics:double-layer}{طول ديباي} الأقصر يحصر التنافر في الفجوات الصغيرة، حيث يغلب
التجاذب، الذي ينمو وفق $1/h$.
\textbf{12.} يتغير \omterm{def:b3:colloids:stability}{تركيز التخثر الحرج} وفق $z^{-6}$ لشحنة الأيون المضاد.
\textbf{13.} $50/64 = \qty{0.78}{mM}$.
\textbf{14.} $50/729 = \qty{0.069}{mM}$.
\textbf{15.} الكاتيونات هي الأيونات المضادة للجسيمات السالبة: فهي
تتراكم في الطبقة المزدوجة وتحجبها.
\textbf{16.} الكبريتات أيون مرافق، تتنافر معه الجسيمات.
\textbf{17.} \qty{0.069}{mol} من \ce{Al^{3+}} لكل متر مكعب.
\textbf{18.} \qty{0.034}{mol} من $\ce{Al2(SO4)3}$.
\textbf{19.} $0.0343 \times 342.3 = \qty{11.7}{g}$.
\textbf{20.} \qty{11.7}{kg} في الساعة.
\textbf{21.} $\ce{Al^3+ + 3 HCO3- -> Al(OH)3 + 3 CO2}$.
\textbf{22.} $3 \times 0.069 = \qty{0.21}{mM}$ من \qty{1.0}{mM}: أي الخُمس. والهيدروجين كربونات الباقية
و\ce{CO2} المذاب يحفظان pH الماء: فلا ينخفض pH إلا
قليلًا.
\textbf{23.} تُمتز أنواع الألومنيوم العالية الشحنة وتعكس شحنة
الجسيمات، فتعود إلى التنافر.
\textbf{24.} \textbf{نحو \qty{12}{g} من كبريتات الألومنيوم لكل متر مكعب} (\qty{11.7}{g})، قبل الزيادة
اللازمة عمليًا للحلمأة.
\end{solution}
